\documentclass[a4paper,12pt]{ctexart}
\usepackage{xcolor}
\usepackage{graphicx}
\usepackage[top=1.8cm, bottom=1.8cm, left=1.8cm, right=1.8cm]{geometry}
\usepackage{array}
\usepackage{hyperref}
\usepackage{booktabs}
\hypersetup{colorlinks=true,linkcolor=blue!70!black}
\linespread{1.35}

\definecolor{storyblue}{RGB}{0,102,153}
\definecolor{thinkorange}{RGB}{230,120,20}
\definecolor{funboxbg}{RGB}{245,248,252}

\newenvironment{thinkbox}{
  \vspace{0.3cm}
  \begin{center}
  \begin{minipage}{0.88\textwidth}
  \colorbox{funboxbg}{
  \begin{minipage}{0.94\textwidth}
  \vspace{0.15cm}
  \textbf{\textcolor{thinkorange}{【 想一想 】}}
}{
  \vspace{0.15cm}
  \end{minipage}}
  \end{minipage}
  \end{center}
  \vspace{0.3cm}
}

\newenvironment{funfact}{
  \vspace{0.2cm}
  \begin{center}
  \begin{minipage}{0.88\textwidth}
  \fbox{
  \begin{minipage}{0.92\textwidth}
  \vspace{0.1cm}
  \textbf{\textcolor{storyblue}{【 有趣的小知识 】}}
}{
  \vspace{0.1cm}
  \end{minipage}}
  \end{minipage}
  \end{center}
  \vspace{0.2cm}
}

\newenvironment{figplaceholder}[1]{
  \vspace{0.2cm}
  \begin{center}
  \fbox{
  \begin{minipage}{0.82\textwidth}
  \centering
  \textcolor{blue!70!black}{\textbf{【插图位置】}} \\[0.2cm]
  #1
  \end{minipage}}
  \end{center}
  \vspace{0.2cm}
}

\title{\textcolor{blue!70!black}{\Huge\textbf{金融学小故事}}}
\author{}
\date{}

\begin{document}
\maketitle
\thispagestyle{empty}

\textcolor{gray}{\small 本模块讲的是《经济小故事》没覆盖的金融知识（股票、债券、基金、保险、汇率等）。建议先读《经济小故事》第1-2、6课了解钱的来历和银行怎么工作。}

\newpage

\section{\textcolor{orange!80!black}{第一课\quad 股票是怎么发明的？——从一条船到一座交易所}}

\subsection*{一条船和很多人的主意}

四百多年前，荷兰的阿姆斯特丹是一个热闹的港口城市。商人们盯上了一个赚钱的好生意——去遥远的东印度群岛（今天的印度尼西亚）运香料回来。

香料在欧洲贵如黄金。一趟成功的航行能赚几十倍的利润。

但是，去东印度的路远得吓人。海上风暴、坏血病、海盗——十艘船出去，能回来七八艘就算运气好了。来回一趟要两三年。

如果你是一个商人，你要面对一个两难的选择：如果自己造一艘船出海，船沉了——你倾家荡产。如果不造船，又赚不到钱。

这时候有人想出了一个绝妙的主意：\textbf{不要一个人承担所有风险，让很多人一起来！}

他把一趟航程需要的钱分成1000份，每一份叫"一股"。你买10股，隔壁王叔叔买50股，对面面包店老板买5股，码头上扛包的工人攒了半年买2股。每个人出的钱不多，但凑在一起，船就造好了。

船出发了。如果平安回来，赚的钱按股分。如果沉了，每个人只亏自己出的小份，谁也不会活不下去。

\begin{figplaceholder}
此处插入"航海集资"插图——一艘帆船周围很多人拿着钱，箭头指向船，标注"风险分摊"
\end{figplaceholder}

这就是\textbf{股份公司}的发明！

\subsection*{从桥头到交易所}

这些"股"买到手后，有人想卖、有人想买。大家慢慢聚在阿姆斯特丹的一座桥上进行交易。后来人太多，桥站不下了，就搬到了附近的一栋楼里。

1602年，荷兰东印度公司正式成立，发行了世界上第一批股票。那栋交易股票的楼，就是世界上第一家证券交易所。

你今天在手机上看到的"上证指数""纳斯达克""恒生指数"，所有那些跳动的数字和红绿箭头，起点就是四百多年前阿姆斯特丹的一座桥。

\begin{thinkbox}
如果你是当时的一个荷兰人，你会买东印度公司的股票吗？最大的好处是什么？最大的风险是什么？如果这条航线跑了十年都赚钱，大家会怎么做？——对了，会有更多人造船、更多人买船票、每条船的利润会被分摊得越来越薄。这就是市场。
\end{thinkbox}

\begin{funfact}
荷兰东印度公司曾经是世界上最大的公司，巅峰时期有超过4万名员工、1万多艘船、15000场战役。但它后来还是破产了。所以买股票不但要选好公司，还要在合适的时候卖出去。
\end{funfact}

\newpage

\section{\textcolor{orange!80!black}{第二课\quad 三个袋子——一个关于花钱的简单计划}}

小明每个月有50块零花钱。以前他的习惯是：拿到钱就花掉，花光拉倒。到了月底想买个东西——没钱了。

这个月他决定换一种方式。他找妈妈要了三个信封，在上面画了标记：

\begin{center}
\begin{tabular}{lll}
\toprule
\textbf{袋子} & \textbf{每月放多少} & \textbf{用来干什么} \\
\midrule
红袋子：零花 & 20块 & 随时花——零食、贴纸、小玩具 \\
蓝袋子：存钱 & 20块 & 存起来以后买大件——比如乐高 \\
绿袋子：投资 & 10块 & 试试能不能让钱变多 \\
\bottomrule
\end{tabular}
\end{center}

\begin{figplaceholder}
此处插入"三个信封"插图——红色(花)、蓝色(存)、绿色(投)，每个里面有模拟纸币
\end{figplaceholder}

\subsection*{半年后发生了什么？}

\textbf{红袋子：}每月花光，0块——但小明吃到了零食、买到了喜欢的贴纸。

\textbf{蓝袋子：}20×6=120块——够买一套乐高忍者城了。

\textbf{绿袋子：}每月10块投进去，小明做了几件事：

\begin{itemize}
  \item 买了同学手工做的书签，拿到学校加价卖——30块进货，卖了52块（+22）
  \item 借给爸爸说好还的时候多5块——60块变65块（+5）
  \item 买了一个热门玩具，玩腻了二手卖掉——30块变22块（-8）
\end{itemize}

总计：绿袋子从90块变成了139块——赚了49块。

\begin{thinkbox}
小明的三个袋子计划合理吗？如果你是小明，你会怎么分配？为什么投资既有赚也有赔？如果小明把所有零花钱都放进绿袋子做"投资"，结果会怎样？
\end{thinkbox}

\subsection*{最重要的三个道理}

\begin{enumerate}
  \item \textbf{钱要分着管}——全花掉、全存着、全拿去冒险，都不如分开放聪明
  \item \textbf{让钱为你工作}——钱不只是用来花的，它可以帮你赚更多的钱
  \item \textbf{有赚就有赔}——投资不是每次都赚钱，所以不要把所有的钱都拿去做高风险的事
\end{enumerate}

\begin{funfact}
世界上最会投资的人之一沃伦·巴菲特，11岁买了人生第一只股票，今年90多岁了还在投资。他从小养成的习惯就是——不花不必要的钱，让每一分钱都持续产生价值。
\end{funfact}

\newpage

\section{\textcolor{orange!80!black}{第三课\quad 保险——为什么每人出一点就能救一家人？}}

\subsection*{小明的自行车被偷了}

小明攒了好久的零花钱，买了一辆心爱的自行车。结果骑了不到一个月——停在楼下被偷了。

小明哭了一场。他爸妈安慰他说："以后再买一辆吧。"可是那要再攒大半年。

后来小明发现，他的同学小红的自行车也是新的，但小红告诉他：她爸爸给她买了一年的自行车盗抢险，只要30块钱。如果自行车被偷了，保险公司赔她一辆新的。

"30块钱就赔你一辆新车？"小明不敢相信。

\subsection*{道理其实很简单}

我们设想一下：一个小区里有1000辆自行车，每年大概有10辆被偷。一辆新车500块。

如果让每一个车主都赔偿被偷的人——那太不公平了，被偷的人太倒霉了。

但如果\textbf{每个人都出一点点钱}放到一个共同的"池子"里：每人出5块，就有5000块。10个人的自行车被偷了，每人拿500块去买新车。每个人的损失是5块，而不是500块。

这就是\textbf{保险}的核心思想：\textbf{很多人各出一点小钱，凑起来帮那些倒霉的人承担大损失。}

但这个"池子"需要精算师来算——每年大概多少人出事、赔多少钱、收多少保费才能刚好覆盖还不亏本。这个数学叫\textbf{精算学}。

\begin{center}
\begin{tabular}{ll}
\toprule
\textbf{没有保险} & \textbf{有保险} \\
\midrule
被偷的人损失500块 & 每人交5块，被偷的人拿500块赔偿 \\
没被偷的人不受影响 & 没被偷的人"损失"5块保费 \\
\textbf{结论}：被偷的人太倒霉了 & \textbf{结论}：小钱化解了大风险 \\
\bottomrule
\end{tabular}
\end{center}

\begin{figplaceholder}
此处插入"保险池"插图——1000个人每人往池子里投硬币，一个自行车被偷的人从池子里拿走一辆新车
\end{figplaceholder}

\subsection*{保险的种类}

不只是自行车才需要保险。大人的保险更多：

\begin{itemize}
  \item \textbf{车险}——开车撞了，保险公司赔修车费和医药费。你不买交强险不能上路
  \item \textbf{医疗保险}——生病住院，保险公司付大部分医药费
  \item \textbf{房屋保险}——如果房子着火了，保险公司赔你一套新的
  \item \textbf{人寿保险}——如果家里的顶梁柱去世了，保险公司给家里一笔钱，让家人能继续生活
\end{itemize}

还有一种很特别的——\textbf{相互保险}。没有保险公司赚差价，所有参保人既是投保人又是保险人。比如农业相互保险——农民们自己组织起来，谁家的庄稼被冰雹打了，大家凑钱帮他。

\begin{thinkbox}
你觉得"保险"这个名字起得好吗？"保"什么？"险"又是什么？如果保险公司的精算师算错了，收的保费太低，会发生什么？
\end{thinkbox}

\begin{funfact}
最早的保险可以追溯到4000年前的中国和巴比伦。中国商人在长江上运货时，会把货物分散装在好几条船上——如果一条船沉了，损失由所有商人一起分担。这就是"分散风险"的雏形。后来犹太商人和意大利商人在此基础上发明了正式的海上保险契约。
\end{funfact}

\newpage

\section{\textcolor{orange!80!black}{第四课\quad 债券是什么？——借钱给国家是什么体验}}

\subsection*{小明想借钱给爸爸}

小明的蓝袋子里存了120块钱。有一天他对爸爸说："爸爸，我把这120块借给你吧。你半年后还我125块，多出来的5块就当利息。"

爸爸笑了："好啊，我给你写个借条。"他拿出一张纸，写上：

\begin{quote}
\textbf{借条} \\
今向小明借款人民币120元整。约定半年后归还，利息5元。\\
借款人：大明（签名）\\
日期：2026年5月23日
\end{quote}

这就是一张最简单的\textbf{债券}！

\begin{figplaceholder}
此处插入"借条"插图——一张手写的借条，上面写借款人、金额、利息、期限
\end{figplaceholder}

\subsection*{从借条到债券}

如果小明爸爸是政府，这张借条就叫\textbf{国债}。如果小明爸爸是一家大公司，这张借条就叫\textbf{公司债}。

国债被认为是最安全的投资之一——因为政府不太可能不还钱（当然也有例外，比如希腊政府2009年就还不上钱了，引发了欧洲债务危机）。

债券有几个关键要素：

\begin{center}
\begin{tabular}{ll}
\toprule
\textbf{要素} & \textbf{小明借条里的对应} \\
\midrule
面值（本金） & 120块 \\
票面利率 & 5/120≈8.3\%（年化） \\
期限 & 半年 \\
发行人 & 爸爸（大明） \\
\bottomrule
\end{tabular}
\end{center}

\begin{thinkbox}
如果爸爸的公司快倒闭了，你还愿意把钱借给他吗？如果你不乐意了，想把手里的借条卖掉——别人愿意出多少钱买？会比120块高还是低？
\end{thinkbox}

\subsection*{债券的价格会波动}

假设小明爸爸的公司突然经营出了问题，大家觉得它可能还不上钱了。这时小明想卖掉手里的借条——别人只愿意出100块买它（因为风险变高了）。这就是\textbf{债券价格下跌}。

反过来，如果公司经营非常好，大家抢着想买它的借条——小明手里的借条可能卖到130块。这就是\textbf{债券价格上涨}。

所以债券虽然看起来比股票"稳"，但它也会涨跌——只是幅度通常比股票小。

\begin{funfact}
2024年，中国10年期国债收益率大约2.3\%。而2008年金融危机前，这个数字是4.5\%左右。为什么现在低了这么多？因为市场利率整体下降了。利率和债券价格是反向关系——利率越低，已发行的债券越值钱。
\end{funfact}

\newpage

\section{\textcolor{orange!80!black}{第五课\quad 基金——让专业的人帮你管钱}}

\subsection*{小明的困惑}

小明想去买股票，但他遇到了几个问题：

\begin{itemize}
  \item 他只有200块钱——买一手（100股）茅台要16万，根本不够
  \item 他不知道买哪家公司的股票——几千家公司，选哪个？
  \item 他也没时间天天盯着股票
\end{itemize}

这时候妈妈告诉他："你可以买基金。"

\subsection*{基金是什么？}

基金就是一个大池子——千千万万的人把钱放进池子里，由一个专业的人（基金经理）帮大家决定买什么股票、买多少。

\begin{center}
\begin{tabular}{lll}
\toprule
\textbf{自己买股票} & \textbf{买基金} \\
\midrule
自己选公司 & 基金经理帮你选 \\
最少买100股 & 10块就能买 \\
得一直盯着 & 不用管 \\
赚了全归你 & 赚了按份分 \\
亏了全你扛 & 亏了大家分摊 \\
\bottomrule
\end{tabular}
\end{center}

\begin{figplaceholder}
此处插入"基金池"插图——很多人往一个大池子里投钱，池子里有各种股票图标，一个戴着眼镜的基金经理在管理
\end{figplaceholder}

\subsection*{基金的种类}

\textbf{按投资对象分：}

\begin{itemize}
  \item \textbf{股票基金}——大部分钱买股票。收益高、风险大
  \item \textbf{债券基金}——大部分钱买债券。收益稳、风险小
  \item \textbf{货币基金}——只买短期、超安全的资产（比如余额宝）。基本不亏，但收益也最低
  \item \textbf{混合基金}——股票债券都买。风险和收益在中间
\end{itemize}

\textbf{按管理方式分：}

\begin{itemize}
  \item \textbf{主动基金}——基金经理自己选股，试图跑赢市场。管理费比较贵
  \item \textbf{指数基金（被动基金）}——不选股，完全复制某个指数（比如沪深300、标普500）。管理费很便宜。巴菲特多次说过：普通投资者最好的选择就是买指数基金
\end{itemize}

\begin{thinkbox}
为什么巴菲特建议普通人买指数基金而不是自己选股？原因有两条：第一，大多数主动基金经理长期跑不赢指数；第二，指数基金管理费低得多。几十年下来，管理费的差距就是一大笔钱。
\end{thinkbox}

\subsection*{策略：定投}

定投就是"定期投入固定金额买基金"。不管基金贵还是便宜，每月固定时间买。

\begin{itemize}
  \item 基金贵的时候，同样的钱买到的份额少
  \item 基金便宜的时候，同样的钱买到的份额多
  \item 长期下来，买入的平均成本被"摊平"了
\end{itemize}

\begin{thinkbox}
假设一只基金第一个月10块/份，第二个月跌到5块/份，第三个月涨回10块/份。如果你每月定投100块——三个月总共投入300块，赚了还是亏了？算算看。答案：第一个月买10份、第二个月买20份、第三个月买10份，共40份，价值400块——赚了100块。这就是定投的"微笑曲线"。
\end{thinkbox}

\begin{funfact}
股神巴菲特在2007年跟一家对冲基金打了一个十年的赌：标普500指数基金的收益会超过任何一只精选的对冲基金组合。十年后，标普500涨了125\%，而那只对冲基金组合只涨了36\%。巴菲特赢了。
\end{funfact}

\newpage

\section{\textcolor{orange!80!black}{第六课\quad 通货膨胀——钱为什么会"变毛"？}}

\subsection*{爷爷的回忆}

小明的爷爷说："我年轻的时候，一根冰棍才五分钱。"
小明问："那现在一根冰棍要多少钱？"
"最便宜也要两块钱了。"

五分钱变成两块钱——涨了40倍。这就是\textbf{通货膨胀}。

\subsection*{通货膨胀是什么？}

简单说：通货膨胀就是\textbf{钱不值钱了}。同样的钱，能买到的东西变少了。

\textbf{为什么会发生？}

最核心的原因：\textbf{钱印得太多了}。

如果社会上东西的总量没变，但钱的总量翻了一倍——那每块钱能买到的东西就少了一半。价格就涨了一倍。

\begin{center}
\begin{tabular}{ll}
\toprule
\textbf{钱的总量} & \textbf{一袋大米的价格} \\
\midrule
100亿 & 100块 \\
200亿 & 200块 \\
500亿 & 500块 \\
\bottomrule
\end{tabular}
\end{center}

\textbf{什么东西会影响钱的总量？} 央行多印钱、银行多放贷、政府发钱刺激经济——都会让市场上流通的钱变多。

\begin{figplaceholder}
此处插入"通货膨胀"对比插图——左边一个购物筐装得满满（20块能买很多），右边同一个筐空荡荡（20块只能买一点）
\end{figplaceholder}

\subsection*{通胀率多少算正常？}

\begin{itemize}
  \item \textbf{0-2\%：}温和通胀。经济学家认为这是"健康"的——人们会倾向于今天花钱而不是等（因为明天更贵），经济有活力
  \item \textbf{2-5\%：}通胀偏高。你的存款购买力在明显缩水
  \item \textbf{5\%以上：}高通胀。生活成本快速上涨
  \item \textbf{10\%以上：}恶性通胀。比如委内瑞拉（最高达到1000000\%），一个月的工资不够买一顿饭
  \item \textbf{负值：}通货紧缩。东西越来越便宜——听起来很好，但大家都不花钱了（因为明天更便宜），工厂倒闭、工人失业
\end{itemize}

\begin{thinkbox}
\textbf{通胀对谁最不公平？}存钱的人。如果你的10000块存银行年利率1.5\%，而通胀率是3\%——你的实际购买力每年缩水1.5\%。反过来说，借钱的人受益——因为10年后还的钱，实际价值已经大大缩水了。这就是为什么通胀本质上是一种隐形的财富再分配。
\end{thinkbox}

\subsection*{北京的通货膨胀}

过去20年，北京什么东西涨得最疯？房价。2005年北京平均房价约6000元/平米，2025年约6万元/平米——涨了10倍。但电子产品的价格反而降了——一台2005年的笔记本电脑卖1万块，性能还不如现在2000块的手机。不同的东西涨价速度不一样，所以通胀率是一个\"一篮子商品\"的平均值——不是你个人的感受。

\begin{funfact}
历史上最严重的恶性通胀发生在2008年的津巴布韦。100万亿津巴布韦元——面值上有14个0——只能买一个面包。后来津巴布韦干脆放弃了本国货币，改用美元。
\end{funfact}

\newpage

\section{\textcolor{orange!80!black}{第七课\quad 金融海啸是怎么回事？——当信任崩塌的时候}}

\subsection*{一个关于信任的游戏}

全班50个同学，每人放10块钱到讲台上的一个盒子里。老师说："盒子里的钱大家可以随时拿回去，但如果等到下个月再拿，每人多给5块。"

一开始没人拿。第二周，有一个人拿回去了——拿到15块。其他人更放心了。第三周，又有几个人拿回去了——都拿到了多出来的钱。

突然有一天传出消息：老师说其实盒子里已经没钱了。所有人冲到讲台前要拿回自己的钱——但盒子里只剩下原来的一半。大部分人只拿回了不到5块。

\textbf{什么都没变，只是"信任"变了。}

\subsection*{2008年的故事}

2008年，美国发生了一件类似的事——规模大了几万倍。

之前几年，美国房价一路上涨。银行想："既然大家都要买房，为什么不贷更多的款？"他们甚至把钱借给了没有稳定收入的人（次贷）。

银行把这些贷款合同打包成一种叫做MBS（住房抵押贷款支持证券）的产品，卖给了全世界的投资者。评级机构给这些产品打了AAA的最高安全评级。

然后房价涨不动了——开始跌。很多借了钱的人发现：房子还没贷款多，不还了。于是大量贷款违约。那些被打了AAA评级的产品——突然变得一文不值。

结果是：拥有大量这类产品的雷曼兄弟银行倒闭了（2008年9月15日）。全球金融市场冻结。股市暴跌。美国失业率从5\%飙升到10\%。全世界陷入了自1930年代大萧条以来最严重的经济危机。

\begin{figplaceholder}
此处插入"2008金融海啸"时间线漫画——从贷款→打包→卖→房价跌→违约→连锁反应
\end{figplaceholder}

\subsection*{教训是什么？}

\begin{enumerate}
  \item \textbf{信任是金融的基石}——一旦信任崩塌，再复杂的金融体系也瞬间瓦解
  \item \textbf{复杂的产品不一定更安全}——看不懂的东西，可能隐藏着巨大的风险
  \item \textbf{高杠杆等于高风险}——银行用很少的本金撬动了巨大的投资，赚的时候暴赚，亏的时候扛不住
\end{enumerate}

\begin{thinkbox}
你有没有经历过类似的事——比如班里传"XXX便利店要关门了"，然后大家全跑去花掉卡里的余额？这就叫"挤兑"。银行挤兑也是一样的道理——本来银行没事，但如果大家全去取钱，银行真的会没钱可给。
\end{thinkbox}

\begin{funfact}
2008年金融危机之后，各国央行总共"印刷"了超过10万亿美元的货币来救市。10万亿美元是什么概念？如果把它换成100美元的钞票叠起来，可以叠到月球那么高。
\end{funfact}

\newpage

\section{\textcolor{orange!80!black}{第八课\quad 汇率——为什么出国旅游有时贵有时便宜？}}

\subsection*{小明的美国梦}

小明的暑假愿望是去美国玩一趟。妈妈算了算说："现在去有点贵。"

为什么？因为\textbf{汇率}变了。

汇率就是一种货币换成另一种货币的价格。如果1美元=7块人民币，那1000美元就需要7000块人民币。如果1美元=6块人民币，那1000美元只需要6000块人民币。

\begin{figplaceholder}
此处插入"汇率天平"插图——天平左边是美元符号，右边是人民币符号，中间指针左右摆动
\end{figplaceholder}

\subsection*{汇率为什么波动？}

\textbf{原因一：国家经济的强弱}

如果一个国家的经济增长快、企业赚钱多、大家对这个国家有信心——它的货币就更值钱。

\textbf{原因二：利率高低}

如果一个国家的利率高，全世界的投资者都把钱存到这个国家来赚利息——大家都来买这个国家的货币，它的汇率就上升。

\textbf{原因三：贸易顺差/逆差}

中国卖很多商品到美国（玩具、手机、家具），美国人付美元——但美元最终要换成人民币才能在中国花。所以市场上有很多人拿着美元要换人民币——人民币的需求高，就会升值。

\begin{center}
\begin{tabular}{ll}
\toprule
\textbf{因素} & \textbf{对人民币汇率的影响} \\
\midrule
中国经济增长快 & 升值 \\
中国央行降息 & 贬值 \\
中美贸易战升级 & 贬值 \\
外国投资者大量买入中国股票 & 升值 \\
\bottomrule
\end{tabular}
\end{center}

\subsection*{汇率对你意味着什么？}

\begin{itemize}
  \item \textbf{如果你要去美国旅游——}汇率高了，你花更多人民币才能换同样多的美元
  \item \textbf{如果你要买进口的东西——}苹果手机、进口奶粉、国外游戏——汇率高了它们更贵
  \item \textbf{如果你的家人从国外汇钱回来——}汇率高了，同样的美元能换更多人民币
  \item \textbf{如果你将来想去留学——}汇率变化可能让你的学费差好几万
\end{itemize}

\begin{thinkbox}
\textbf{探究题：}去查一下过去5年人民币对美元的汇率走势。2021年是什么水平？2024年是什么水平？找一找中间发生了什么大事让你看到的那些变化。
\end{thinkbox}

\begin{funfact}
2022年，因为美联储大幅加息，美元指数涨到了20年来的最高点。人民币从6.3贬值到了7.3左右——如果你当时有1万美元要换成人民币，同一个人、同一笔钱，一年多出了1万块人民币。
\end{funfact}

\end{document}
